\newcommand{\DoNotLoadEpstopdf}{}
\documentclass[11pt]{article}
\usepackage[T1]{fontenc}
\usepackage{lmodern}
\usepackage[margin=1in]{geometry}
\usepackage{amsmath,amssymb,amsthm,mathtools}
\usepackage{microtype}
\usepackage{enumitem}
\usepackage{booktabs,array}
\usepackage{xcolor}
\usepackage{tikz}
\usetikzlibrary{matrix,fit,positioning,backgrounds}
\usepackage{needspace}
\usepackage[hidelinks,unicode]{hyperref}
\input{glyphtounicode}
\input{glyphtounicode-cmex}
\pdfgentounicode=1
\pdftrailerid{}
\hypersetup{pdftitle={Pointwise Multiple Ergodic Averages for Mixing Transformations},pdfauthor={OpenAI}}
\newtheorem{theorem}{Theorem}[section]
\newtheorem{proposition}[theorem]{Proposition}
\newtheorem{lemma}[theorem]{Lemma}
\newtheorem{corollary}[theorem]{Corollary}
\theoremstyle{definition}
\newtheorem{definition}[theorem]{Definition}
\theoremstyle{remark}
\newtheorem{remark}[theorem]{Remark}
\newcommand{\R}{\mathbb R}
\newcommand{\Z}{\mathbb Z}
\newcommand{\N}{\mathbb N}
\newcommand{\C}{\mathbb C}
\newcommand{\T}{\mathbb T}
\newcommand{\E}{\mathbb E}
\newcommand{\dd}{\,d}
\newcommand{\law}{\mathcal L}
\newcommand{\norm}[1]{\lVert#1\rVert}
\newcommand{\abs}[1]{\lvert#1\rvert}
\newcommand{\ip}[2]{\langle#1,#2\rangle}
\DeclareMathOperator{\tr}{tr}
\DeclareMathOperator{\Ran}{Ran}
\DeclareMathOperator{\supp}{supp}
\DeclareMathOperator{\rank}{rank}
\DeclareMathOperator{\spanop}{span}
\newcommand{\articleid}[1]{\begingroup\hyphenchar\font=-1\relax\textnormal{#1}\endgroup}
\setlist[enumerate]{itemsep=3pt,topsep=5pt}
\setlist[itemize]{itemsep=3pt,topsep=5pt}
\numberwithin{equation}{section}
\title{Pointwise Multiple Ergodic Averages\\for Mixing Transformations}
\author{OpenAI}
\date{October 4, 2026}
\begin{document}
\maketitle
\begin{abstract}
Let $T$ be an invertible mixing probability-preserving transformation.
For every integer $n\ge2$ and every fixed tuple of bounded measurable
functions, we prove that the consecutive multiple ergodic averages of
length $n$ converge almost everywhere to the product of the integrals,
as the averaging length tends to infinity through all positive integers.
The probability space need not be standard, and no rate of mixing is
required.
\end{abstract}
\section{Introduction}
\label{sec:introduction}

Let $(X,\mathcal F,\mu)$ be a probability space and let $T:X\to X$
be invertible, bimeasurable, and measure preserving. We say that $T$
is \emph{mixing} if
\begin{equation}\label{eq:mixing}
 \mu(A\cap T^{-r}B)\longrightarrow\mu(A)\mu(B)
 \qquad (|r|\to\infty)
\end{equation}
for every $A,B\in\mathcal F$. For a fixed tuple of bounded measurable
functions, the consecutive multiple ergodic averages are
\[
 A_N(f_1,\ldots,f_n)(x)
 =\frac1N\sum_{k=1}^N\prod_{j=1}^n f_j(T^{jk}x).
\]
The problem is to determine their almost-everywhere behavior, not merely
their limit in norm. We prove the following result.

\begin{theorem}\label{thm:main}
Suppose that $T$ satisfies \eqref{eq:mixing}. For every integer
$n\ge2$ and every fixed $f_1,\ldots,f_n\in L^\infty(\mu)$,
\[
 A_N(f_1,\ldots,f_n)(x)
 \longrightarrow\prod_{j=1}^n\int_X f_j\dd\mu
 \qquad\text{for $\mu$-almost every }x,
\]
as $N$ tends to infinity through all positive integers. The probability
space need not be standard, and no rate of mixing is required.
\end{theorem}

The exceptional null set may depend on the system and the fixed tuple.
The theorem concerns the times $k,2k,\ldots,nk$ and does not require
a common exceptional set for all bounded functions.

\subsection{Context and prior work}

Arithmetic-progression averages became central to ergodic Ramsey theory
through Furstenberg's proof of Szemer\'edi's theorem
\cite{Furstenberg1977}. The distinction between norm and pointwise
convergence is substantial. Host and Kra proved $L^2$ convergence of
these averages for every finite length on arbitrary
probability-preserving systems \cite[Theorem~1.1]{HostKra2005}.
Ziegler obtained another proof through universal characteristic
factors \cite[Corollary~1.8]{Ziegler2007}.
For pointwise convergence, Bourgain's double recurrence theorem treats
two factors \cite{Bourgain1990}. Higher-length results have been
obtained under additional structural assumptions: K-systems
\cite[Theorem~4.2]{DerrienLesigne1996}, weakly mixing systems with
singular spectrum on their Pinsker factor \cite{Assani1998}, and
weakly mixing systems whose pairwise independent self-joinings are
independent \cite[Theorem~3.4]{GutmanHuangShaoYe2018}. Huang, Shao,
and Ye proved pointwise convergence for every finite length on
ergodic measure-distal systems \cite[Theorem~C]{HuangShaoYe2019}.

The present proof extends the fourfold mixing theorem of
\cite[Theorem~1.1]{FourfoldCompanion}. It uses two substantial
companion results. The first proves that ordinary mixing implies
mixing of every finite order, the assertion of Rokhlin's
multiple-mixing problem \cite[Theorem~1.1]{MixingCompanion}.
The second gives an oscillation estimate for triple averages at
distinct integer slopes on arbitrary invertible probability-preserving
systems \cite[Proposition~3.5]{SlopesCompanion}; together with its
profile approximation, this yields pointwise triple convergence
at distinct positive integer slopes without a mixing assumption.
Section~\ref{sec:inputs} states their exact forms and derives the
required triple pointwise theorem along all lengths. These companion
proofs are not reproduced here. The selected-law construction,
arbitrary-length extension argument, and tensor estimate used below
are proved in full.

\subsection{The arbitrary-length step}

We briefly describe what must change when the number of factors
increases. A failure of pointwise convergence allows an averaging
length to be chosen measurably at each starting point so that a
centered average stays positive. After passing to a suspension flow,
these choices define a probability law on $n$ output coordinates.
The law has the correct individual marginals and is invariant under
different speeds in the different coordinates. The same construction
works on extensions of the flow.

The method of \cite[Sections~3--6]{FourfoldCompanion} organizes
additional outputs into three groups. Their spatial coordinates are
independent, and the original spatial coordinate is independent of
any two groups; independence from all three groups is not imposed.
We retain this three-group structure even when there are $n$ slots.
Higher-order mixing and a tensor-packing estimate show that these
independence properties survive the selected joining construction.

Two changes make this approach work for arbitrary $n$. First, we
minimize the number of slots supporting a nonzero centered correlation.
This gives independence for every smaller collection of the specific
spatial lists being tested. A slot may read both the original spatial
coordinate and the outputs in any proper subset of the three groups.
That flexibility permits the outputs to be regrouped without imposing
independence on arbitrary collections of whole states.

Second, an extension is chosen so that conditioning on all distal
parameters together with a single state cannot improve under any
further extension. This is a conditional-expectation energy
construction, related to the sated-extension method
\cite{Austin2015} and adapted from
\cite[Section~6]{FourfoldCompanion}. Applying the joining construction
again gives a square array. The single-state conditioning identity
controls its distinguished row and column. Their parameter laws and
a Hilbert-space argument then produce a new nonzero correlation with
the same number of slots but one more output-only slot. Maximizing
that number gives the contradiction. Only the pointwise theorem for
at most three slots on arbitrary systems is used.

Section~\ref{sec:inputs} records the dynamical inputs.
Section~\ref{sec:selected} constructs the selected laws and their
conditional marginals. Section~\ref{sec:channels} proves preservation
of the three-group structure. Sections~\ref{sec:correlations}
and~\ref{sec:square} carry out the minimal-correlation and square
arguments. Appendix~\ref{sec:packing} proves the tensor estimate.

\section{Dynamical inputs and standard models}
\label{sec:inputs}

Fix an integer $n\ge2$ and put $I=\{1,\ldots,n\}$. We first state
the dynamical results used below. The distinction between norm and
pointwise convergence matters: norm convergence is available for any
finite number of slots, whereas the pointwise input on arbitrary
systems will be used only for at most three slots.

\subsection{Standard models and the two companion inputs}

For the fixed tuple in Theorem~\ref{thm:main}, record all integer
translates of bounded measurable representatives of the functions.
This gives a measurable map into a countable product of compact scalar
discs, intertwining $T$ with the invertible shift. The pushforward
measure defines a standard Borel probability system, and mixing passes
to this factor. The original functions are its coordinate functions.
Changing representatives affects at most a countable union of translated
null sets, so almost-everywhere convergence on the factor pulls back.
We may therefore work on standard probability spaces.

Throughout the proof, factors and conditional laws are understood
modulo null sets. We use Borel versions of measurable functions.
Sub-$\sigma$-algebras on a standard probability space have standard
models modulo null sets. A flow is a jointly Borel action
$(S_Y^t)_{t\in\R}$ on a standard Borel probability space $(Y,m_Y)$,
preserving $m_Y$. Factor maps commute with every fixed time modulo
null sets. Joint measurability permits Fubini in all time integrals.
For a preserved factor, induced transformations at fixed times can be
modeled as invertible maps; we take time integrals on the original flow.

We use the following statement from
\cite[Theorem~1.1]{MixingCompanion}.
\begin{theorem}[Higher-order mixing, companion input]
\label{thm:multiple-mixing}
Let $T$ be an invertible mixing probability-preserving transformation.
For any fixed $s\ge2$, $g_1,\ldots,g_s\in L^\infty(\mu)$,
and integer times $t_1,\ldots,t_s$,
\[
 \int_X\prod_{j=1}^s g_j(T^{t_j}x)\dd\mu(x)
 \longrightarrow\prod_{j=1}^s\int_Xg_j\dd\mu
 \quad\text{as}\quad
 \min_{i\ne j}|t_i-t_j|\longrightarrow\infty.
\]
For $s=1$ the equality holds for every time.
\end{theorem}
The cited theorem is stated for sets and ordered times. Simple-function
approximation, invariance, and the finitely many orderings of the times
give this formulation. No quantitative mixing rate is asserted or used.

Our second companion input is the smooth-average convergence consequence
of \cite[Proposition~3.5 and Lemma~3.6]{SlopesCompanion}.
An \emph{averaging profile} there is a real compactly supported smooth
function $\varphi$ of mass one with moments of orders $1,2,3$ equal to
zero, and, for some $G\ge1$,
\[
 \norm{\varphi^{(r)}}_\infty
 \le \bigl(G(r+2)\bigr)^{G(r+2)}\qquad(r\ge0).
\]
The constant $G$ may depend on the profile. We need the following
two precise conclusions.
\begin{proposition}[Smooth triple averages, companion input]
\label{prop:smooth-input}
Let $R$ be any invertible probability-preserving transformation, and
let $b_1,b_2,b_3$ be distinct nonzero integers with greatest common
divisor one. For bounded $g_1,g_2,g_3$ and every averaging profile,
\[
 V_k^\varphi(y)=\sum_{l\in\Z}2^{-k}\varphi(l/2^k)
              \prod_{j=1}^3g_j(R^{b_jl}y)
\]
converges almost everywhere as $k\to\infty$. Moreover, for each
$c\in[1,2]$ and $\epsilon>0$ there is such a profile with
\[
 \norm{\varphi-c^{-1}\mathbf1_{(0,c]}}_{L^1(\R)}<\epsilon.
\]
\end{proposition}
The approximation is \cite[Lemma~3.9]{SlopesCompanion}. Neither part
requires mixing or ergodicity. We will not assume pointwise convergence
of longer averages on arbitrary systems.

\subsection{Relative weak mixing and distal factors}

We recall the structure facts in the form needed for the auxiliary
systems. An extension of an invertible system $R$ over a factor $B$
is \emph{relatively weakly mixing} if, for bounded $f,g$ with
$\E[f\mid B]=0$,
\begin{equation}\label{eq:rwm}
 \frac1H\sum_{h=1}^H
 \norm{\E[\overline f(g\circ R^h)\mid B]}_2^2\longrightarrow0.
\end{equation}
An extension is \emph{relatively compact} if the invariant finitely
generated $L^\infty(B)$-modules span a dense subspace of its $L^2$
space. A module consists of finite sums $\sum_j a_jF_j$, with
$a_j\in L^\infty(B)$ and fixed $F_j\in L^2$; invariance is under
composition with both $R$ and $R^{-1}$. Individual modules need not be
closed; see the corrected module characterization in
\cite[Theorem~2.5(ii)$'$]{JamneshanCorrigendum2026}.
A \emph{relatively distal} extension is a tower of relatively
compact extensions and inverse limits. A system is \emph{distal} if
it is distal over the trivial factor. A joining is an invariant
probability law on a product with the prescribed marginals.

\begin{proposition}[Relative structure]\label{prop:structure}
The following statements hold for standard probability-preserving
systems, without an ergodicity assumption.
\begin{enumerate}
\item Every system $Y$ has a maximal measure-distal factor $W_Y$,
and $Y\to W_Y$ is relatively weakly mixing. Every distal factor
pulls into the maximal distal factor of an extension. Automorphisms
commuting with the transformation preserve $W_Y$.
\item Factors and joinings of distal systems are distal. A distal
system is relatively distal over each of its factors. Distality is
equivalent for a transformation and any of its nonzero integer powers.
\item Relative weak mixing passes to nonzero integer powers. The
conditional product of relatively weakly mixing fibre laws over an
invariant joining of their base laws is relatively weakly mixing
over that joining.
\item A relatively weakly mixing extension and a relatively distal
extension of the same base are relatively disjoint: every invariant
joining identifying their bases is the conditional product over
that base.
\item For a distal transformation $R$, every $d\ge1$, and bounded
$g_1,\ldots,g_d$, the averages
\[
 \frac1N\sum_{l=1}^N\prod_{i=1}^d g_i\circ R^{il}
\]
converge almost everywhere.
\end{enumerate}
For a flow, the maximal distal factor of time one is preserved by
the flow, and its action at every nonzero rational time is distal.
\end{proposition}

These are the conclusions of
\cite[Proposition~2.2]{FourfoldCompanion}. They include the
nonergodic form of the Furstenberg--Zimmer structure theorem
\cite{Zimmer1976,Jamneshan2023} and the pointwise theorem for
distal systems \cite[Theorem~C]{HuangShaoYe2019}. The cited proposition
gives the component-disintegration argument needed for the latter
in the nonergodic setting. In the flow assertion, invariance of
$W_Y$ follows because each time commutes with time one. If
$r=p/q\ne0$ is rational, the $q$th power of its action is the
distal action at integer time $p$, so the power assertion applies.

The following short argument, from
\cite[Lemma~2.3]{FourfoldCompanion}, will also fix the norm-convergence
input explicitly.
\begin{lemma}[Relative norm convergence]\label{lem:relative-norm}
Suppose $Y\to B$ is relatively weakly mixing for $R$. For distinct
nonzero integers $a_1,\ldots,a_d$ and bounded $G_1,\ldots,G_d$,
\begin{equation}\label{eq:relative-norm}
 \left\|\frac1N\sum_{l=1}^N
 \left(\prod_{i=1}^dG_i\circ R^{a_il}
       -\prod_{i=1}^d\E[G_i\mid B]\circ R^{a_il}\right)\right\|_2
 \longrightarrow0.
\end{equation}
\end{lemma}
\begin{proof}
By telescoping and scaling, it suffices to show that the average of
$v_l=\prod_i G_i\circ R^{a_il}$ tends to zero when all functions
are bounded by one and one of them, $G_s$, is centered over $B$.
Induct on $d$. For fixed $h\ge1$ put
$F_{i,h}=\overline{G_i}(G_i\circ R^{a_ih})$. Invariance gives
\[
 \ip{v_l}{v_{l+h}}
 =\int_Y F_{s,h}\prod_{i\ne s}
                  F_{i,h}\circ R^{(a_i-a_s)l}\dd m_Y.
\]
The induction hypothesis replaces the factors indexed by $i\ne s$
in their average by their conditional expectations over $B$. The
remaining product is $B$-measurable and bounded by one; hence
\[
 \limsup_{N\to\infty}
 \left|\frac1N\sum_{l=1}^N\ip{v_l}{v_{l+h}}\right|
 \le\norm{\E[F_{s,h}\mid B]}_1.
\]
For $d=1$ this bound holds directly. The Hilbert-space van der
Corput inequality now bounds the limiting squared norm by a constant
times
\[
 \frac1H+\frac1H\sum_{h=1}^H\norm{\E[F_{s,h}\mid B]}_1.
\]
Relative weak mixing for $R^{a_s}$ and Cauchy--Schwarz make this
expression tend to zero as $H\to\infty$.
\end{proof}

For positive slopes, inserting functions $1$ in omitted positions
and using Proposition~\ref{prop:structure}(5) shows that the
distal-factor averages have an $L^2$ limit. Thus
Lemma~\ref{lem:relative-norm} gives an $L^2$ limit for arbitrary
bounded tuples with distinct positive slopes on every standard
system. On the original mixing system, the trivial factor already
satisfies \eqref{eq:rwm}; the limit is the product of the means.

\subsection{Triple convergence along all lengths}

\begin{lemma}[Triple pointwise convergence]\label{lem:triple-convergence}
On any standard invertible probability-preserving system, averages
of at most three bounded functions at distinct positive integer
slopes converge almost everywhere along all positive integer lengths.
\end{lemma}
\begin{proof}
Add constant-one functions if necessary to have three distinct
positive slopes, write them as $gb_1,gb_2,gb_3$ with
$\gcd(b_1,b_2,b_3)=1$, and apply
Proposition~\ref{prop:smooth-input} to $S=R^g$. By scaling, suppose
the three functions are bounded by one, and put
$F_l=\prod_{j=1}^3g_j\circ S^{b_jl}$.

For each rational $c\in[1,2]$, choose a sequence of profiles tending
in $L^1$ to $h_c=c^{-1}\mathbf1_{(0,c]}$. On a common conull set
all the corresponding smooth dyadic averages converge. For
$N_k=\lfloor c2^k\rfloor$, the difference between the ordinary
average at $N_k$ and $V_k^\varphi$ is at most
\[
 \frac1{c2^k}
 +2^{-k}\sum_{l\in\Z}|h_c(l/2^k)-\varphi(l/2^k)|.
\]
The last sum tends to $\norm{h_c-\varphi}_1$ by Riemann sums.
Arbitrarily close profiles therefore give the Cauchy property on
each of these scales. The $L^2$ limit established above identifies
all their almost-everywhere limits with the same function.

For any bounded sequence $|F_l|\le1$, changing the averaging length
from $N$ to $M$ changes the average by at most
\begin{equation}\label{eq:length-change}
                 \frac{2|M-N|}{\max(M,N)}.
\end{equation}
For $2^k\le N<2^{k+1}$, compare $N$ with
$\lfloor c2^k\rfloor$ on a finite rational grid in $[1,2]$.
Equation~\eqref{eq:length-change} makes the limiting error
arbitrarily small as the grid mesh tends to zero. This proves
convergence along all integer lengths.
\end{proof}

\section{Selected lengths and their joining laws}\label{sec:selected}

We assume that the asserted pointwise convergence fails for the fixed
set of slopes $I=\{1,\ldots,n\}$.  This section records that failure in
a probability law and constructs compatible laws on all flow extensions.
The construction retains a positive correlation while forcing every
one-, two-, or three-slot marginal to have its ordinary conditional law.
We follow the selector and joining construction in
\cite[Sections~3--6]{FourfoldCompanion}, keeping the extension and
conditioning properties explicit for the arbitrary number of slots.

\subsection{A positive witness at selected lengths}

Expanding each function into its mean and centered part, and then into
real and imaginary parts, gives a nonempty $J_0\subseteq I$ and real
functions $f_i$, $i\in J_0$, such that
\[
  \int_X f_i\,\dd\mu=0,\qquad \norm{f_i}_\infty\le1,
  \qquad
  B_L(x)=\frac1L\sum_{l=1}^L\prod_{i\in J_0}f_i(T^{il}x)
\]
does not tend to zero on a set of positive measure.  Here and below
$B_L$ has an integer subscript.  Lemma~\ref{lem:relative-norm}, applied
over the trivial factor, gives $B_L\to0$ in $L^2(\mu)$.  After changing
the sign of one $f_i$, there are $\delta>0$ and a measurable set
$E_*\subseteq X$ of positive measure such that
\begin{equation}\label{eq:failure-witness}
  B_L(x)>2\delta\quad\text{for arbitrarily large integers }L,
  \qquad x\in E_*.
\end{equation}

We use the roof-one suspension
\begin{equation}\label{eq:suspension}
  C=X\times[0,1),\qquad m_C=\mu\otimes\dd u,
  \qquad
  S_C^t(x,u)=\bigl(T^{\lfloor u+t\rfloor}x,\{u+t\}\bigr).
\end{equation}
Write $\xi(x,u)=x$ and $\upsilon(x,u)=u$, regarding $\upsilon$ also as
a coordinate in the circle $\R/\Z$.

\begin{lemma}[Selected positive witness]\label{lem:selected-witness}
There are a finite set $F\subseteq[1,2]$ containing $1$ and measurable
functions $N_a:C\to[2^a,\infty)$, $a\in\N$, independent of the phase,
with
\begin{equation}\label{eq:selector-grid}
  N_a(x,u)\in\{c2^k:c\in F,\ k\ge a\},\qquad
  N_a(x,u)=2^a\quad(x\notin E_*),
\end{equation}
such that $B_{\lfloor N_a(x,u)\rfloor}(x)>\delta$ on $E_*$.
Moreover, there is a continuous nonnegative function $\Psi$ on
$(\R/\Z)^I$ for which
\begin{equation}\label{eq:selected-positive-limit}
 \liminf_{a\to\infty}\int_C\frac1{N_a(z)}
 \int_0^{N_a(z)}
 \Psi\bigl((\upsilon(S_C^{it}z))_{i\in I}\bigr)
 \prod_{i\in J_0}f_i\bigl(\xi(S_C^{it}z)\bigr)
 \,\dd t\,\dd m_C(z)>0.
\end{equation}
\end{lemma}

\begin{proof}
By~\eqref{eq:length-change}, changing the length from $L$ to $M$
changes an average of a sequence bounded by one by at most
$2\abs{M-L}/\max(M,L)$.  Choose $F$ to be a sufficiently fine finite
grid in $[1,2]$.  For every sufficiently large integer $L$, some
$c\in F$ and $k=\lfloor\log_2L\rfloor$ satisfy
\[
  \abs{B_L(x)-B_{\lfloor c2^k\rfloor}(x)}<\delta
  \qquad\text{for every }x.
\]
Equation~\eqref{eq:failure-witness} therefore supplies, for each
$x\in E_*$ and $a$, a pair $(c,k)$ with $k\ge a$ at which the latter
average exceeds $\delta$.  Choose the first such pair in a fixed
enumeration.  This is a measurable choice from a countable set and
defines $N_a$ on $E_*$; use $2^a$ elsewhere.

Choose continuous nonnegative circle functions $\alpha,\beta$, neither
identically zero, supported respectively inside $(1/4,1/2)$ and
$(0,1/(4n))$.  Set
\begin{equation}\label{eq:phase-witness}
 \Psi(v)=\alpha(2v_1-v_2)\beta(v_2-v_1),
 \qquad v=(v_i)_{i\in I}\in(\R/\Z)^I.
\end{equation}
For $z=(x,u)$ and $t=l+s$, where $l\in\Z$ and $0\le s<1$, the
phase factor is $\alpha(u)\beta(s)$.  If it is nonzero, then
$0<u+is<1$ for every $i\in I$, so
$\xi(S_C^{it}z)=T^{il}x$.  Summing complete unit intervals gives,
uniformly in $z$ and real $L\ge2$,
\begin{align}
 &\frac1L\int_0^L
 \Psi\bigl((\upsilon(S_C^{it}z))_i\bigr)
 \prod_{i\in J_0}f_i\bigl(\xi(S_C^{it}z)\bigr)\,\dd t
 \notag\\
 &\hspace{18mm}
 =\alpha(u)\left(\int_0^1\beta(s)\,\dd s\right)
 B_{\lfloor L\rfloor}(x)+O(L^{-1}).
 \label{eq:suspension-witness}
\end{align}
The error includes the incomplete last interval, the change from
indices $0,\ldots,\lfloor L\rfloor-1$ to $1,\ldots,\lfloor L\rfloor$,
and the normalization.  On $E_*$ the main term at $L=N_a(z)$ has
integral at least
\[
 \delta\mu(E_*)
 \left(\int_0^1\alpha(u)\,\dd u\right)
 \left(\int_0^1\beta(s)\,\dd s\right)>0.
\]
On $X\setminus E_*$ its integral tends to zero because $N_a=2^a$
there and $B_{2^a}\to0$ in $L^2(\mu)$.  The uniform error is
$O(2^{-a})$, proving~\eqref{eq:selected-positive-limit}.
\end{proof}

\subsection{Selected laws on all extensions}

Fix these selectors and a free ultrafilter $\mathcal U$ on $\N$.
A \emph{flow over $C$} is a standard probability space $(Y,m_Y)$
with a jointly Borel probability-preserving action $S_Y^t$, $t\in\R$,
and a probability-preserving factor map $\pi_Y:Y\to C$ satisfying
$\pi_Y S_Y^t=S_C^t\pi_Y$ modulo null sets for each fixed $t$.
An extension $\rho:Y'\to Y$ of such flows is a probability-preserving
factor map, equivariant in the same sense, with
$\pi_{Y'}=\pi_Y\rho$ modulo null sets.  We do not require a common
conull set on which every real-time factor identity holds.  Joint
measurability and Fubini justify their uses inside time integrals.

\begin{proposition}[Selected joining laws]\label{prop:selected-law}
Every flow $Y$ over $C$ has a unique probability law $\lambda_Y$ on
$Y^I$ satisfying, for all bounded measurable functions $G_i$ on $Y$,
\begin{equation}\label{eq:selected-law}
 \int_{Y^I}\prod_{i\in I}G_i(y_i)\,\dd\lambda_Y
 =\lim_{a\to\mathcal U}\int_Y\frac1{N_a(\pi_Yy)}
   \int_0^{N_a(\pi_Yy)}\prod_{i\in I}G_i(S_Y^{it}y)
   \,\dd t\,\dd m_Y(y).
\end{equation}
Its marginals are $m_Y$, and it is invariant under
$(y_i)_i\mapsto(S_Y^{it}y_i)_i$, $t\in\R$.
If $\rho:Y'\to Y$ is an extension over $C$, then
$(\rho^I)_*\lambda_{Y'}=\lambda_Y$.

Countable inverse limits of extensions over $C$ exist in this
category.  The selected law at such a limit projects to the selected
law at every stage.  Pullbacks of the stage selected-law spaces have
dense union in $L^2$ of the limit selected law.
\end{proposition}

\begin{proof}
For each $a$ the expression before the ultralimit defines a genuine
probability measure $\lambda_{Y,a}$ on $Y^I$, by integrating the
measurable orbit map over the indicated finite time interval.  We
first show that each of its marginals converges against every bounded
measurable test $G$ to $m_Y$.  The one-slot continuous-time averages
\[
 \frac1L\int_0^L G(S_Y^{it}y)\,\dd t
\]
converge almost surely as $L\to\infty$.  For example, apply Birkhoff's
theorem for $S_Y^i$ to the function
$y\mapsto\int_0^1G(S_Y^{is}y)\,\dd s$ and discard the incomplete
unit interval.  Their almost-sure limit has integral $\int G\,\dd m_Y$,
since every deterministic average has that integral.  As
$N_a(\pi_Yy)\ge2^a$, bounded convergence gives the same integral
limit at the selected lengths.

Embed $Y$ as a Borel subset of a compact metrizable space $K$.
Ultralimits of the $\lambda_{Y,a}$ against continuous functions on
$K^I$ give a positive normalized functional, hence a probability
measure on $K^I$.  Its marginals are $m_Y$, by the preceding paragraph,
so it is concentrated on $Y^I$.  Call the resulting law $\lambda_Y$.
To obtain~\eqref{eq:selected-law} for arbitrary bounded slot functions,
normalize their bounds to one and approximate each in $L^1(m_Y)$ by
continuous functions on $K$ with the same bound.  The difference
between the two product integrals is bounded by the sum of the
single-slot errors.  Under $\lambda_Y$ these errors are the
$L^1(m_Y)$ errors, while under $\lambda_{Y,a}$ they converge to those
errors by the marginal convergence already proved.  Letting the
errors tend to zero proves the formula.  Product tests determine a
probability measure on $Y^I$, giving uniqueness.  In particular, all
these identities depend only on the almost-everywhere classes of
their slot functions.

Applying the product action at a fixed time $t_0$ shifts the time
interval in~\eqref{eq:selected-law} by $t_0$ without changing its
length or its initial-point selector.  For slot functions bounded by
one the error is at most $2\abs{t_0}2^{-a}$, proving invariance.
For an extension $\rho$, substitute $G_i\circ\rho$ in the formula.
Equivariance, the equality of the selectors, and Fubini identify its
right-hand side with that for $Y$, proving functoriality.

For a tower $Y_0\leftarrow Y_1\leftarrow\cdots$, take the full
countable product of the underlying standard spaces, with the
consistent inverse-limit probability measure and coordinatewise
flow.  This is a jointly Borel action.  A cylinder through stage $r$
has probability computed from $m_{Y_r}$ by the finitely many bonding
maps.  At any fixed time their equivariance and preservation of
$m_{Y_r}$ leave that probability unchanged.  Cylinders determine the
measure, so every fixed time preserves it.  The measure is
concentrated on compatible sequences; using
the full product avoids any need for one compatibility set invariant
at all real times.  The stage projections are factor maps modulo
null sets, and projection through $Y_0$ gives the factor onto $C$.
Functoriality now gives the asserted selected-law projections.
Finally, the stage coordinates generate the inverse-limit
sigma-algebra.  The compatibility relations hold under the selected
law slot by slot because its marginals are the inverse-limit
probability measure.  Thus the pulled-back stage sigma-algebras on
the selected-law space are increasing and generate it modulo null
sets.  The usual $L^2$ approximation by an increasing sequence of
sigma-algebras proves density.
\end{proof}

\subsection{Distal parameters and small marginals}

For a flow $Y$ over $C$, let
$\omega_Y:Y\to W_Y$ be its maximal measure-distal factor for time one,
and let $\nu_Y=(\omega_Y)_*m_Y$.  Write $m_{Y,w}$ for its conditional
probability measures, so that
\[
 m_Y=\int_{W_Y}m_{Y,w}\,\dd\nu_Y(w).
\]
Under $\lambda_Y$ put
\begin{equation}\label{eq:distal-parameters}
 P_i=\omega_Y(y_i),\qquad P=(P_i)_{i\in I}.
\end{equation}
Proposition~\ref{prop:structure} ensures that $W_Y$ is preserved by
every flow time.  Such statements about factor sigma-algebras are
understood modulo null sets; all time integrals below are evaluated
on the original jointly Borel flow.  If $\rho:Y'\to Y$ is an extension,
the pullback of $W_Y$ is distal for time one and hence is contained
in $W_{Y'}$.  Consequently the old parameter array is a measurable
function of the new parameter array on selected-law spaces.

\begin{proposition}[Parameter and small-marginal laws]
\label{prop:selected-parameters}
For every flow $Y$ over $C$ and bounded measurable $H_i$ on $W_Y$,
\begin{equation}\label{eq:deterministic-parameters}
 \E_{\lambda_Y}\prod_{i\in I}H_i(P_i)
 =\lim_{L\to\infty}\int_Y\frac1L\int_0^L
   \prod_{i\in I}H_i\bigl(\omega_Y(S_Y^{it}y)\bigr)
   \,\dd t\,\dd m_Y(y).
\end{equation}
For every nonempty $J\subseteq I$ with $\abs{J}\le3$,
\begin{equation}\label{eq:small-conditional-law}
 \law_{\lambda_Y}\bigl((y_i)_{i\in J}\mid P\bigr)
   =\bigotimes_{i\in J}m_{Y,P_i}.
\end{equation}
\end{proposition}

\begin{proof}
We first explain the passage from discrete to continuous time used
in both assertions.  Write $t=l+s$, with $l\ge0$ an integer and
$0\le s<1$.  For each fixed $s$, a continuous-time product becomes
the discrete product at slopes $i$ for the shifted slot functions
$G_i\circ S_Y^{is}$.  Whenever the corresponding discrete averages
converge almost surely, Fubini and bounded convergence show that
their integrals in $s$ converge almost surely.  Starting the discrete
average at $l=0$ rather than $l=1$, and dropping an incomplete final
unit interval, each has an error tending to zero.  The resulting
continuous-time convergence holds along all real lengths $L$.

Apply this observation to functions from $W_Y$.  Their shifts still
belong to $W_Y$, and pointwise convergence on the distal factor is
part of Proposition~\ref{prop:structure}.  Thus the time averages
in~\eqref{eq:deterministic-parameters} converge almost surely along
all lengths, and therefore at the selected lengths as well.  Bounded
convergence and~\eqref{eq:selected-law} prove that equation.

For $\abs{J}\le3$, Lemma~\ref{lem:triple-convergence} and the same
continuous-time argument show that the selected $J$-marginal equals
its deterministic-length integrated law.  The extension
$Y\to W_Y$ is relatively weakly mixing.  Apply
Lemma~\ref{lem:relative-norm} to the distinct integer slopes
$i\in J$ and the shifted functions $G_i\circ S_Y^{is}$, for each
fixed $s$.  Conditional expectation commutes with these shifts
because $W_Y$ is preserved.  Integration in $s$, using the uniform
norm bound and bounded convergence, gives
\begin{equation}\label{eq:small-replacement}
 \E_{\lambda_Y}\prod_{i\in J}G_i(y_i)
 =\E_{\lambda_Y}\prod_{i\in J}
    \E_{m_Y}[G_i\mid W_Y](y_i).
\end{equation}
Indeed the difference of the deterministic averages tends to zero
in $L^2(m_Y)$, so also in their integrated values.  Applying this
identity to $G_i$ multiplied by arbitrary bounded functions of
$P_i$ yields
\begin{equation}\label{eq:small-own-parameters}
 \law_{\lambda_Y}\bigl((y_i)_{i\in J}\mid(P_i)_{i\in J}\bigr)
 =\bigotimes_{i\in J}m_{Y,P_i}.
\end{equation}

Conditioning on the full parameter array requires one more step;
it does not follow from~\eqref{eq:small-own-parameters} alone.
Use the invariant transformation with speed $i$ in slot $i$.
By the stability under nonzero powers and conditional products in
Proposition~\ref{prop:structure}, the $J$-marginal system
in~\eqref{eq:small-own-parameters} is relatively weakly mixing over
its parameter law.  The full $P$-system is a joining of the distal
systems $W_Y$ at these integer powers.  It is distal, and is
therefore relatively distal over the subarray $(P_i)_{i\in J}$.
Their joint law under $\lambda_Y$ is an invariant joining
identifying that common subarray.  Relative disjointness in
Proposition~\ref{prop:structure} makes this joining conditionally
independent over the subarray, which is precisely
\eqref{eq:small-conditional-law}.
\end{proof}

The full parameter law is thus unaffected by the selectors, although
the full state law need not be.  We now turn that state law into a new
flow while preserving one chosen slot as the original system.

\begin{lemma}[Normalized joining flow]\label{lem:normalized-joining}
For $b\in I$, define
\begin{equation}\label{eq:normalized-joining}
 \mathsf M_b(Y)=
 \left(Y^I,\lambda_Y,
   S_{\mathsf M_b(Y)}^t(y_i)_i=(S_Y^{it/b}y_i)_i\right).
\end{equation}
This is a standard probability-preserving flow over $C$ via
$(y_i)_i\mapsto\pi_Y(y_b)$, and projection onto slot $b$ is an
extension onto $Y$.  The array $P$ belongs to the maximal
measure-distal factor of $\mathsf M_b(Y)$ for time one.
\end{lemma}

\begin{proof}
Invariance follows from Proposition~\ref{prop:selected-law}, after
rescaling the time parameter by $1/b$.  Slot $b$ has speed one and
marginal $m_Y$, giving the stated factor maps.  In slot $i$ the
$W_Y$-factor is distal for time $i/b$, by
Proposition~\ref{prop:structure}.  The law of $P$ is an invariant
joining of these rational-time distal systems, hence is distal.
Maximality of the distal factor gives the last assertion.
\end{proof}

\begin{lemma}[The retained positive correlation]\label{lem:positive-witness}
With the functions and phase test of Lemma~\ref{lem:selected-witness},
\begin{equation}\label{eq:positive-witness}
 \int_{C^I}\Psi\bigl((\upsilon(z_i))_i\bigr)
    \prod_{i\in J_0}f_i\bigl(\xi(z_i)\bigr)\,\dd\lambda_C>0.
\end{equation}
The phase array in this integral is measurable in the distal
parameter array of $C^I$.
\end{lemma}

\begin{proof}
Continuous functions on the compact product of circles are uniformly
approximable by finite sums of products of continuous one-coordinate
functions.  Apply this to $\Psi$, use~\eqref{eq:selected-law} for
each resulting product, and then use the strictly positive lower
limit in~\eqref{eq:selected-positive-limit}.  Finally, time one fixes
the phase on $C$, so the phase factor is distal and lies in $W_C$.
\end{proof}

\section{Channels and closure under selected laws}\label{sec:channels}

The selected length is determined by the original point in the suspension.
It can therefore change the joint law of other coordinates, even when
their unconditioned law is a product.  We organize those coordinates in
three groups.  Independence of the original point from any two groups
will control the effect of selecting a length on all three groups.
This channel construction and its closure argument adapt
\cite[Section~5]{FourfoldCompanion}; the proof below permits every fixed
finite set of slots $I=\{1,\ldots,n\}$.

\begin{definition}[Channel]\label{def:channel}
A \emph{channel} is a flow $Y$ over the suspension $C$, with factor map
$\pi_Y:Y\to C$, together with a finite list of probability-preserving
maps $D_\gamma:Y\to C$, called \emph{outputs}, and positive rational
numbers $q_\gamma$, such that
\[
 D_\gamma(S_Y^t y)=S_C^{q_\gamma t}D_\gamma(y)
\]
for each fixed $t$, modulo null sets.  Partition the output indices into
three groups, numbered $1,2,3$; empty groups are allowed.  Set
\[
 x_0=\xi\circ\pi_Y,\qquad x_\gamma=\xi\circ D_\gamma,
 \qquad d_Y=(x_\gamma)_\gamma,
\]
and let $d_{Y,H}$ be the sublist from groups in
$H\subseteq\{1,2,3\}$.  Let $V$ collect the phase coordinates of
$\pi_Y$ and all the $D_\gamma$.  Under the measure $m_Y$ require
\begin{equation}\label{eq:channel-laws}
 \begin{split}
  \law(d_Y\mid V)&=\bigotimes_\gamma\mu,\\
  \law(x_0,d_{Y,H}\mid V)&=
       \mu\otimes\bigotimes_{\gamma\text{ in groups }H}\mu,
       \qquad H\subsetneq\{1,2,3\}.
 \end{split}
\end{equation}
An empty spatial list has the one-point law.
\end{definition}

Thus the spatial outputs have a product law independent of the phases.
The original spatial coordinate has a product law with the outputs in
any proper collection of groups, but no independence from the complete
output list is required.  The suspension $C$, with no outputs, is a
channel.

\begin{lemma}[Conditioning and extensions]\label{lem:channel-conditioning}
For a channel $Y$, both identities in \eqref{eq:channel-laws} hold with
conditioning on its maximal measure-distal factor $W_Y$ in place of
$V$.  Every flow extension of $Y$ over $C$ is a channel with the pulled
back output data, and the same strengthening holds there.
\end{lemma}

\begin{proof}
Choose a positive integer $h$ such that $hq_\gamma$ is an integer for
each output.  The phase vector $V$ is a factor fixed by time $h$, and
its time-one action is distal, so $V$ belongs to $W_Y$.  For any spatial
tuple appearing in \eqref{eq:channel-laws}, its joint factor with $V$
has product law over $V$.  At time $h$ the spatial coordinates evolve
by a product of positive integer powers of $T$.  This product is mixing:
the assertion follows first for product tests from mixing of each
power, and then for all bounded tests by approximation.  Consequently
this factor is relatively weakly mixing over $V$ for time $h$.

On the other hand, $W_Y$ is distal for time $h$ and relatively distal
over $V$.  Relative disjointness, in the form of
\cite[Proposition~2.2]{FourfoldCompanion}, makes the spatial factor and
$W_Y$ relatively independent over $V$.  Its conditional spatial law
given $W_Y$ is therefore the product law in
\eqref{eq:channel-laws}.  Pullback to an extension preserves the joint
laws of the original data, proving the channel identities there; the
same argument gives their stronger conditioning.
\end{proof}

We now prove the closure property that will permit repeated extensions.
Recall from Lemma~\ref{lem:normalized-joining} that
$\mathsf M_b(Y)$ has state space $Y^I$, measure $\lambda_Y$, and flow
$(y_i)_i\mapsto(S_Y^{it/b}y_i)_i$.

\begin{proposition}[Channel closure]\label{prop:channel-closure}
Let $Y$ be a channel and $b\in I$.  On $\mathsf M_b(Y)$ take
$\pi_Y(y_b)$ as original factor and take every $D_\gamma(y_i)$,
$i\in I$, as an output, in the old group of $\gamma$.
These data make $\mathsf M_b(Y)$ a channel.
\end{proposition}

\begin{proof}
Each new output is probability-preserving and has positive rational
speed $iq_\gamma/b$.  It remains to prove the two conditional laws.
Test them with products of bounded spatial functions and products of
bounded phase functions.  Such tests determine the laws.  Expanding
each spatial function into its mean and centered part, and scaling,
reduces the proof to real spatial functions that are either $1$ or
centered and bounded in absolute value by $1$.  We may also bound every
phase function by $1$.  Whenever at least one spatial factor is
centered, the tested integral must be zero.  A test using the new
original coordinate uses outputs from only a proper set of groups.

\paragraph{Time intervals and offsets.}
Use the selected-law formula \eqref{eq:selected-law} and condition on
the old phase vector $V$ under $m_Y$.  Fix an integer $h>0$ clearing
all output-speed denominators and write
$t=hl+s$, where $l\ge0$ is an integer and $0\le s<h$.
For fixed $V,s$, all phase weights in the test are independent of $l$.
The spatial function for output $\gamma$ at slot $i$ is evaluated at
\begin{equation}\label{eq:channel-times}
 T^{e_\gamma i l+a_{\gamma,i}}x_\gamma,
 \qquad e_\gamma=hq_\gamma\in\N.
\end{equation}
The offsets $a_{\gamma,i}$ depend on $V,s$ and range over a fixed
finite set, by the suspension formula.  If the original-factor test
is present, its argument is $T^{hbl+a_0}x_0$, with the same finite-offset
property.  Equivariance and Fubini justify these formulas in the time
integrals.

On an event where the selected length is $N=c2^k$, $c\in F$ and
$k\ge a$, put $m=\lfloor c2^k/h\rfloor$.  Replacing the normalized
time integral by the average over the $m$ complete intervals, followed
by integration against $ds/h$, has error at most $2h2^{-a}$.
For all sufficiently large $a$ we have $m\ge1$.

\paragraph{Tests using a proper set of groups.}
Given $V$, the relevant old output coordinates have product law and
are independent of $x_0$, by \eqref{eq:channel-laws}.  The selector is
determined by $V,x_0$.  Suppose first that some output occurrence is
centered.  For fixed $V,s,x_0$, integrate the spatial output product
at index $l$.  The expectation factors over the old output indices
$\gamma$.  For an index having a centered occurrence, the times
\eqref{eq:channel-times} in distinct slots separate as $l\to\infty$.
Mixing of all orders makes its integral tend to the product of the
means, hence to zero.  A single centered occurrence already has zero
integral.  The convergence is uniform over the finitely many offset
patterns.  All other factors have absolute value at most $1$.
The Ces\`aro averages of these absolute expectations consequently tend
to zero uniformly over all sufficiently large lengths.  This applies
to the selected lengths as well.

If every output test is $1$, the centered factor must be the original
one.  Birkhoff's theorem for the ergodic power $T^{hb}$ gives convergence
to zero for $\mu$-almost every $x_0$, simultaneously for all the
finitely many offsets $a_0$.  Since the selected lengths tend to
infinity pointwise, they have the same limit.  The conditional law of
$x_0$ given $V$ is $\mu$, so bounded convergence completes this case.

\paragraph{Three groups and almost orthogonal vectors.}
It remains to treat the full output list, with no original-factor
spatial test.  If one group has no centered occurrence, every spatial
test in that group is $1$, and the preceding argument applies; the
phase tests of that group were allowed throughout.  Thus suppose each
group contains a centered occurrence.  For fixed $V,s$, write its
spatial product at index $l$ as
\[
 u_l\in L^2(\Omega_1),\qquad
 v_l\in L^2(\Omega_2),\qquad
 w_l\in L^2(\Omega_3),
\]
where $\Omega_j$ carries the product measure with one copy of $\mu$
for every old output in group $j$.  These are real unit-ball vectors,
indeed functions bounded by $1$.  Given $V$, the old spatial output
law is the product of these three spaces.  The vectors depend on
$V,s$ only through the finite offset patterns.

For every $\eta>0$, after omitting finitely many initial indices, the
families $(u_l)$ and $(v_l)$ each have absolute pair inner products at
most $\eta$ outside a graph of bounded maximum degree $L_0$.
All these bounds are uniform in the offset patterns.  To verify this,
in each of the two groups choose a centered occurrence at some output
index $\gamma$.  Its coordinate integral in a pair inner product is
small by mixing of all orders whenever the times in the two copies of
\eqref{eq:channel-times} are sufficiently separated.  The times within
each copy separate once $l,l'$ exceed a fixed threshold.  Failure of
separation between copies then requires one of finitely many inequalities
\[
 \abs{e_\gamma i l-e_\gamma i'l'+a'}\le R.
\]
For fixed $l$, each inequality permits only a bounded number of
integers $l'$, because $e_\gamma,i'>0$; the same holds with $l,l'$
interchanged.  Their union is the required bounded-degree graph.
The coordinate integrals for other output indices are bounded by $1$.
The separation threshold and $L_0$ may depend on $\eta$, but no rate
of mixing is needed.

\paragraph{Conditioning on a selected length.}
For fixed selector index $a$, partition into the events
$N_a=c2^k$, $c\in F$, $k\ge a$, choosing one representation when a
length has more than one.  Given $V$, write their probabilities as
$p_{k,c}(V)$.  On one such event of probability $p>0$, the conditional
spatial law of any two groups, and of each single group, is unchanged.
Indeed the event is determined by $V,x_0$, which is independent of any
proper set of groups by \eqref{eq:channel-laws}.  The full three-group
law may change.  Its density $\rho$ relative to the product law satisfies
\begin{equation}\label{eq:channel-density}
 0\le p\rho\le1,\qquad
 \int p\rho=p,\qquad \norm{p\rho}_2\le\sqrt p.
\end{equation}
The reason is that the joint measure of the outputs and the event is
dominated by the unconditioned product measure.  These statements hold
for almost every $V$; the length events are countable.

We choose the rank bound to meet two requirements: the packing
bound must be $o(m)$, and the reciprocal bounds must be summable over
dyadic length scales. For $m=\lfloor c2^k/h\rfloor$ set
\begin{equation}\label{eq:channel-dimension}
 d_m=\left\lceil(1+\log m)^{3/2}\right\rceil,
 \qquad z_l=u_l\otimes v_l\quad(0\le l<m).
\end{equation}
Let $\Pi_m$ be the orthogonal projection onto up to $d_m$ leading
eigenvectors of the positive covariance operator
$\sum_{l<m}z_lz_l^*$, taking its full range if its rank is smaller.
Fix $0<\varepsilon\le1$ and then fix
$0<\eta\le\varepsilon^2/10^4$ in the preceding graph bound.
Theorem~\ref{thm:packing}, proved in Appendix~\ref{sec:packing} from
\cite[Theorem~4.1]{FourfoldCompanion}, gives
\begin{equation}\label{eq:channel-projected-average}
 \frac1m\sum_{l<m}\norm{\Pi_m z_l}
       \le\varepsilon+o(1).
\end{equation}
Here the error tends to zero uniformly over the offset patterns.
Indeed, after the fixed initial indices are removed, the number of
projection norms at least $\varepsilon$ is at most
\[
 C(1+L_0)\varepsilon^{-2}
 \exp\!\bigl(C\varepsilon^{-2}\sqrt{d_m}\log(2+d_m)\bigr)=o(m).
\]
The removed indices contribute $o(1)$ to the average.  Importantly,
$\varepsilon,\eta$, and the resulting exceptional degree are fixed
before $m$ tends to infinity.

For the projected terms, Cauchy--Schwarz under the conditional event
law uses its unchanged two-group and one-group marginals.  Term by
term, including the event probability, it gives
\[
 \left|p\int (\Pi_mz_l)(\omega_1,\omega_2)
                   w_l(\omega_3)\rho(\omega)\,d\omega\right|
 \le p\norm{\Pi_mz_l}_2\norm{w_l}_2
 \le p\norm{\Pi_mz_l}_2.
\]
Thus their average contributes at most $p(\varepsilon+o(1))$.

For the residual vectors $z_l'=(1-\Pi_m)z_l$, let
$A=(\langle z_l',z_{l'}'\rangle)_{l,l'<m}$ and
$B=(\langle w_l,w_{l'}\rangle)_{l,l'<m}$.  The covariance trace is
at most $m$, so its remaining eigenvalues, and hence the nonzero
eigenvalues of $A$, are at most $m/d_m$.  If its rank was at most
$d_m$, the residual is zero.  Also $B$ is positive semidefinite with
trace at most $m$.  Consequently
\[
 \left\|\frac1m\sum_{l<m}z_l'\otimes w_l\right\|^2
   =\frac{\operatorname{tr}(AB)}{m^2}\le\frac1{d_m}.
\]
By \eqref{eq:channel-density}, the residual contribution, including
the event probability, is at most $\sqrt{p/d_m}$.
Multiplication by the fixed phase weight does not increase either bound.

\paragraph{Summing over lengths.}
Write $m(k,c)=\lfloor c2^k/h\rfloor$.  For fixed $h$ this is comparable
to $2^k$ for all sufficiently large $k$, uniformly in the finite set
$F$.  Since $\sum_{k,c}p_{k,c}(V)=1$, Cauchy--Schwarz yields
\begin{equation}\label{eq:channel-residual-sum}
 \sum_{k\ge a,\ c\in F}\sqrt{\frac{p_{k,c}(V)}{d_{m(k,c)}}}
 \le
 \left(\sum_{k\ge a,\ c\in F}\frac1{d_{m(k,c)}}\right)^{1/2}
 \longrightarrow0.
\end{equation}
The series converges because its summands are $O(k^{-3/2})$.
The projected contributions have total limit superior at most
$\varepsilon$, by \eqref{eq:channel-projected-average} and the same
probability identity.  These estimates are uniform in $V,s$ and may
be integrated.  There is no measurable-choice issue for the
projections: choose one for each of the countably many sizes and
finitely many offset patterns.  Letting $a$ tend to infinity in the
selected-law formula and then letting $\varepsilon$ tend to zero
proves the required vanishing.  Both laws in
\eqref{eq:channel-laws} follow.
\end{proof}

\begin{corollary}\label{cor:channel-copied-law}
For every channel $Y$, the copied spatial list
$(d_Y(y_i))_{i\in I}$ under $\lambda_Y$ has its full product law
conditional on $P=(\omega_Y(y_i))_{i\in I}$.
\end{corollary}

\begin{proof}
Apply Proposition~\ref{prop:channel-closure} and
Lemma~\ref{lem:channel-conditioning} to $\mathsf M_b(Y)$.
By Lemma~\ref{lem:normalized-joining}, the old parameter array $P$
belongs to its maximal distal factor.  Conditioning the resulting
constant product law down to $P$ proves the assertion.
\end{proof}

\section{Minimal correlations and saturated extensions}
\label{sec:correlations}

The positive correlation in \eqref{eq:positive-witness} records the
assumed pointwise failure. We continue the channel-and-saturation
argument of \cite[Sections~5--6]{FourfoldCompanion}, with an additional
minimum-support choice. We choose such a correlation with the
fewest possible spatial slots. This choice supplies the product laws
for proper subcollections that a pointwise theorem of greater length
would otherwise have to provide. Among correlations of this minimum
size, we maximize the number of slots that use only outputs. The final
section will increase that number and obtain a contradiction.

\subsection{Spatial roles and minimum support}

Let $Y$ be any channel. At a slot $i\in I$, we allow either of the
following spatial lists:
\begin{align*}
 L_i(y_i)&=d_Y(y_i)
   &&\text{(an output role)},\\
 L_i(y_i)&=(x_0(y_i),d_{Y,H_i}(y_i)),\qquad
   H_i\subsetneq\{1,2,3\}
   &&\text{(an input role)}.
\end{align*}
An input role may therefore read the original coordinate jointly with
outputs from any proper subset of the three groups. This allowance is
important: later, regrouping outputs will produce precisely these
joint lists. A test may ignore some entries of its designated list.

For a nonempty set $J\subseteq I$, designate one role at each $i\in J$.
Choose bounded real functions $h_i$ of the corresponding lists, centered
for their product spatial laws. By Lemma~\ref{lem:channel-conditioning},
$h_i(L_i)$ also has conditional mean zero over $W_Y$. We consider
correlations
\begin{equation}
 \E_{\lambda_Y}\left[\Phi(P)
          \prod_{i\in J}h_i(L_i(y_i))\right],
 \label{eq:role-correlation}
\end{equation}
where $\Phi$ is any bounded real function of the full distal parameter
array $P=(\omega_Y(y_i))_{i\in I}$. The channel $C$, with no outputs,
has a nonzero correlation of this form by
Lemma~\ref{lem:positive-witness}.

Let $k$ be the minimum of $|J|$ among all nonzero correlations
\eqref{eq:role-correlation}, over all channels and all such data.
The small-slot law \eqref{eq:small-conditional-law} gives $k\geq4$:
when $|J|\leq3$, the conditional expectation of the product given $P$
is the product of the conditional means, all zero. Thus $n<4$ already
contradicts the assumed failure. In what follows $4\leq k\leq n$.

\begin{lemma}[Product laws below the minimum support]
\label{lem:minimum-product}
On any channel, choose fewer than $k$ distinct slots and designate
one allowed spatial list at each. Under its selected law, these lists,
jointly, have their product spatial law conditional on the full distal
parameter array.
\end{lemma}

\begin{proof}
Test the conditional law by a product of bounded real functions, one
for each designated list, and by a bounded real function of the full
parameter array. Write each spatial test as its product-law mean plus
a centered function. Every nonconstant term in the expansion is a
correlation of the form \eqref{eq:role-correlation} with fewer than $k$
slots, and hence vanishes. Only the product of the means remains.
Product tests determine the asserted conditional law on these finite
products of standard spaces.
\end{proof}

Among the nonzero correlations with support size $k$, let $r$ be the
largest number of output roles. Corollary~\ref{cor:channel-copied-law}
implies $r<k$, because the full copied output collection has product
spatial law given $P$. Choose a channel and a correlation attaining
these two extremal values. We next pass to an extension on which
conditioning on the parameters, or on the parameters and one slot,
cannot be improved by any further extension.

\subsection{Saturating the parameter and single-slot data}

For any flow $Y$ over $C$, define sub-$\sigma$-algebras of the selected
law space $(Y^I,\lambda_Y)$ by
\begin{equation}
 \mathcal D_0(Y)=\sigma(P),\qquad
 \mathcal D_i(Y)=\sigma(P,y_i)\quad(i\in I).
 \label{eq:data-algebras}
\end{equation}
If $\rho:Y'\to Y$ is a flow extension over $C$, functoriality gives
$(\rho^I)_*\lambda_{Y'}=\lambda_Y$. The pullback of $W_Y$ is a distal
factor of $Y'$, so is contained in $W_{Y'}$. Consequently the pullback
of each $\mathcal D_l(Y)$ is contained in $\mathcal D_l(Y')$.

\begin{lemma}[Simultaneous saturation]
\label{lem:saturation}
Every standard flow $Y$ over $C$ has a standard flow extension $Y_*\to Y$
over $C$ with the following property. For every further standard flow
extension $\rho:Y'\to Y_*$ over $C$, every $f\in L^2(\lambda_{Y_*})$,
and every $l\in\{0\}\cup I$,
\begin{equation}
 \E_{\lambda_{Y'}}[f\circ\rho^I\mid\mathcal D_l(Y')]
 =\E_{\lambda_{Y_*}}[f\mid\mathcal D_l(Y_*)]\circ\rho^I.
 \label{eq:saturation}
\end{equation}
\end{lemma}

\begin{proof}
We use a countable bounded-energy construction of the kind underlying
Austin's sated extensions \cite[Theorem~3.11]{Austin2015}; the argument
below proves the assertion for the particular data
\eqref{eq:data-algebras}. In a tower of flow extensions,
for a fixed test $f$ introduced at one stage, the quantity
\[
  \left\|\E[f\mid\mathcal D_l]\right\|_2^2
\]
at later stages is nondecreasing, by inclusion of the pulled-back
data, and is bounded by $\|f\|_2^2$. At each stage introduce a countable
dense set in the $L^2$ space of its selected law, retaining all
previously introduced tests. Schedule all pairs consisting of a retained
test and a datum index so that each pair is treated infinitely often,
with positive tolerances tending to zero. A diagonal enumeration gives
such a schedule even though new tests are introduced at every stage.

At a scheduled step, extend the current flow so that the corresponding
energy is within the prescribed tolerance of its supremum over all
standard flow extensions of that stage over $C$. The supremum is
finite. It can be taken over a set: standard Borel spaces, their
probability laws, Borel actions, and factor maps admit codes on fixed
standard spaces. The identity extension is among the competitors.
Let $Y_*$ be the countable inverse limit of the resulting tower,
which exists in our category by Proposition~\ref{prop:selected-law}.

Fix a retained test, an index $l$, and a further extension $Y'\to Y_*$.
At every scheduled optimization of this pair, $Y'$ is also an extension
of the stage just before the optimization. Its energy is therefore at
most the energy achieved immediately after that optimization plus its
tolerance. The latter achieved energy is at most the energy at $Y_*$.
As the tolerances tend to zero, the energy at $Y'$ is at most that at
$Y_*$. The reverse inequality follows from monotonicity, so the two
energies are equal.

Identify $L^2(\lambda_{Y_*})$ with its pullback in
$L^2(\lambda_{Y'})$. The two conditional expectations are orthogonal
projections onto nested closed subspaces. Equality of their squared
norms forces equality of the projections, proving
\eqref{eq:saturation} for each retained test. Stage functions have
dense span in $L^2(\lambda_{Y_*})$ by
Proposition~\ref{prop:selected-law}; the retained dense sets and the
contraction property of conditional expectation extend the identity
to every $f$ in that space.
\end{proof}

Apply Lemma~\ref{lem:saturation} to the chosen channel. Pulling back
its data preserves the channel laws by
Lemma~\ref{lem:channel-conditioning}, and preserves its nonzero
correlation by functoriality. Its old parameter test is a function of
the enlarged distal parameters. We henceforth call this saturated
channel $Y$. Write $J$ for its support, $A\subset J$ for its output
roles, and choose
\[
  b\in J\setminus A,
  \qquad |J|=k,\quad |A|=r.
\]
Thus the role at $b$ is an input role. The purpose of the remaining
construction is to turn it into an output role while keeping a
nonzero correlation on the same $k$ slots.

\subsection{Regrouping the other slots as outputs}

Set $Z=\mathsf M_b(Y)$. Its state is $(y_j)_{j\in I}$, its law is
$m_Z=\lambda_Y$, and projection to $y_b$ is a flow factor $Z\to Y$.
We give $Z$ a channel structure different from the full-copy structure
of Proposition~\ref{prop:channel-closure}. Retain all old outputs from
slot $b$ with their original group numbers. For each $j\in J\setminus\{b\}$,
also include the $C$-valued factors whose spatial coordinates form
\[
 Q_j=L_j(y_j).
\]
These are old output factors and, when the role at $j$ is an input
role, the original factor $\pi_Y(y_j)$. Choose a surjection
\[
 \tau:J\setminus\{b\}\longrightarrow\{1,2,3\},
\]
possible because $k\geq4$, and put the entire list $Q_j$ in new group
$\tau(j)$. Every new output has the correct marginal and a positive
rational speed: the old speed at slot $j$ is multiplied by $j/b$.

\begin{lemma}[The regrouped channel]
\label{lem:regrouped-channel}
The preceding output lists and groups make $Z$ a channel over its
original factor $\pi_Y(y_b)$.
\end{lemma}

\begin{proof}
We prove both channel laws conditional on the old full parameter
array $P$ on $(Y^I,\lambda_Y)$. Every phase in the new channel is a
function of $P$, so conditioning down then proves
\eqref{eq:channel-laws} for $Z$.

For the full output collection, its spatial lists are
\[
  d_Y(y_b),\qquad Q_j\quad(j\in J\setminus\{b\}).
\]
These are allowed lists at $k$ distinct old slots. Expand a product
of tests on these lists into constant and centered parts. Every term
with fewer than $k$ centered factors has the product-law value by
Lemma~\ref{lem:minimum-product}. The all-centered term also vanishes
against every bounded parameter test. Otherwise it would give a
nonzero correlation with support $k$ and $r+1$ output roles: slot $b$
now has an output role, and the other roles are unchanged. This
contradicts the definition of $r$. If the output list at $b$ is empty,
its centered test is identically zero, with the same conclusion.
Thus the full output collection has product spatial law given $P$.

For the second channel law, choose a proper subset
$H\subsetneq\{1,2,3\}$ of the new groups and include $x_0(y_b)$.
At slot $b$ the resulting list is
$(x_0(y_b),d_{Y,H}(y_b))$, an allowed input list. At any other retained
slot $j$ it is exactly $Q_j$, and such a slot occurs only when
$\tau(j)\in H$. Surjectivity of $\tau$ omits at least one of the
$k-1$ other slots. Fewer than $k$ old slots therefore occur, so
Lemma~\ref{lem:minimum-product} gives the required product law
conditional on $P$. This proves the second channel law as well.
\end{proof}

\section{The square argument}
\label{sec:square}

We now use the saturated channel $Y$, its nonzero correlation
\eqref{eq:role-correlation}, and the regrouped channel
$Z=\mathsf M_b(Y)$ from Section~\ref{sec:correlations}. The selected
law $\lambda_Z$ is a law on an ordered square of $Y$-states. Its
$b$th row and $b$th column both have law $\lambda_Y$, but for different
reasons. Saturation will relate their correlations without requiring
any symmetry of the square or any pointwise theorem for many slots
on an arbitrary extension.

\subsection{A common corner and two conditional expectations}

On $(Y^I,\lambda_Y)$, put
\begin{equation}
 q((y_i)_{i\in I})=
    \prod_{j\in J\setminus\{b\}}h_j(L_j(y_j)),
 \qquad
 g(P,y_b)=\E_{\lambda_Y}[q\mid P,y_b].
 \label{eq:q-and-g}
\end{equation}
The function $g$ is not zero in $L^2(\lambda_Y)$: otherwise the
nonzero correlation would vanish after conditioning on $(P,y_b)$.
On the same state space regarded as $Z$, the function $q$ reads only
the newly added outputs. Write $e=\omega_Z(z)$ for the new distal
parameter and define
\begin{equation}
 u(e,y_b)=\E_{m_Z}[q\mid e,y_b].
 \label{eq:u}
\end{equation}
All these functions have bounded real Borel versions.

By Lemma~\ref{lem:normalized-joining}, the old array $P$ is a factor
of the new distal parameter. Write
\[
 R(e)=P,
 \qquad \kappa(e)=R(e)_b.
\]
Thus $\kappa\circ\omega_Z$ is the old distal parameter of the
slot-$b$ factor $Z\to Y$.

A point of the selected law space $(Z^I,\lambda_Z)$ has the form
\[
 \mathbf z_i=(y_{ij})_{j\in I},\qquad
 E_i=\omega_Z(\mathbf z_i),\qquad E=(E_i)_{i\in I}.
\]
The first index specifies the row. Define the column states and their
parameters by
\begin{equation}
 y_i=y_{ib},\qquad
 p=(\kappa(E_i))_{i\in I},\qquad
 w=\kappa(E_b).
 \label{eq:square-parameters}
\end{equation}
In particular $y_b=y_{bb}$ is the common corner. The column
$(y_i)_{i\in I}$ has law $\lambda_Y$, by functoriality for $Z\to Y$.
The row $\mathbf z_b$ also has law $\lambda_Y$, because the one-slot
marginal of $\lambda_Z$ is $m_Z=\lambda_Y$.
Figure~\ref{fig:square} records these distinct identifications.

\begin{figure}[tbp]
\centering
\begingroup
\setlength{\arraycolsep}{10pt}
\renewcommand{\arraystretch}{1.45}
\[
\begin{array}{c|ccccc|l}
 &1&\cdots&b&\cdots&n&\text{row parameter}\\ \hline
 1&y_{11}&\cdots&\colorbox{blue!10}{$y_{1b}$}&\cdots&y_{1n}&E_1\\
 \vdots&\vdots&&\colorbox{blue!10}{$\vdots$}&&\vdots&\vdots\\
 b&\colorbox{green!12}{$y_{b1}$}&
   \colorbox{green!12}{$\cdots$}&
   \fcolorbox{black}{yellow!25}{$y_{bb}$}&
   \colorbox{green!12}{$\cdots$}&
   \colorbox{green!12}{$y_{bn}$}&E_b\\
 \vdots&\vdots&&\colorbox{blue!10}{$\vdots$}&&\vdots&\vdots\\
 n&y_{n1}&\cdots&\colorbox{blue!10}{$y_{nb}$}&\cdots&y_{nn}&E_n
\end{array}
\]
\[
 \underbrace{\mathbf z_b=(y_{bj})_{j\in I}}_{\text{row }b:
       \text{ old parameters }R(E_b)}
 \qquad
 \underbrace{(y_{ib})_{i\in I}}_{\text{column }b:
       \text{ old parameters }p}
 \qquad
 w=p_b=R(E_b)_b.
\]
\endgroup
\caption{The ordered square under $\lambda_Z$. The highlighted row
and column meet at $y_{bb}$, whose old distal parameter is $w$.
Both have law $\lambda_Y$: the row by the single-slot marginal,
the column by the factor map $Z\to Y$. No invariance under
transposition is asserted.}
\label{fig:square}
\end{figure}


\begin{lemma}[The row--column identity]
\label{lem:row-column}
Under $\lambda_Z$,
\begin{equation}
 \E\big[q((y_i)_i)q(\mathbf z_b)\mid E\big]
   =\int_Y g(p,y)u(E_b,y)\,\dd m_{Y,w}(y).
 \label{eq:row-column}
\end{equation}
Moreover, in a single $Z$-state with law $m_Z$, the conditional law
of its slot-$b$ state given $e$ is $m_{Y,\kappa(e)}$.
\end{lemma}

\begin{proof}
Apply saturation \eqref{eq:saturation} with datum index $0$ to the
extension $Z\to Y$. It says that the conditional law of the entire
column given $E$ is the law of $(Y^I,\lambda_Y)$ given $p$.
This assertion follows first for $L^2$ tests, and then as an identity
of conditional kernels by a countable determining class. By the
singleton case of \eqref{eq:small-conditional-law}, the conditional
law of $y_b$ given $E$ is consequently $m_{Y,w}$. Conditioning down
to $E_b$ gives the same law for the slot-$b$ state of row $b$ given
$E_b$. Since row $b$ has marginal $m_Z$, this proves the last
assertion of the lemma.

Saturation with datum index $b$ gives the stronger identity
\begin{equation}
 \E_{\lambda_Z}[q((y_i)_i)\mid E,\mathbf z_b]=g(p,y_b).
 \label{eq:column-single-row}
\end{equation}
Indeed the old datum is $(P,y_b)$ and the new datum is
$(E,\mathbf z_b)$. Multiply \eqref{eq:column-single-row} by
$q(\mathbf z_b)$ and condition on $E$. The singleton case of
\eqref{eq:small-conditional-law}, now for $Z$, states that the law
of $\mathbf z_b$ given $E$ is the $m_Z$-fibre law over $E_b$.
Within that fibre, \eqref{eq:u} conditions $q(\mathbf z_b)$ on its
slot-$b$ state, whose law is $m_{Y,w}$ by the first part of the proof.
This gives \eqref{eq:row-column}.
\end{proof}

\subsection{Independence of the two parameter descriptions}

The integral in \eqref{eq:row-column} pairs a column function with
a row function. Its nonvanishing will follow from the fact that,
apart from the parameter at their common corner, the two parameter
descriptions are independent.

\begin{lemma}[Parameter independence]
\label{lem:parameter-independence}
Under $\lambda_Z$, the random variables $p$ and $E_b$ are
conditionally independent given $w$.
\end{lemma}

\begin{proof}
Let $F_0$ be a bounded real test on $W_Z$. In the single-state
parameter law $(\omega_Z)_*m_Z$, let $\overline F_0$ satisfy
\[
 \overline F_0(\kappa(e))=\E[F_0(e)\mid\kappa(e)].
\]
We prove that $F_0(E_b)$ can be replaced by
$\overline F_0(w)$ when tested against any product
$\prod_i H_i(p_i)$ of bounded real coordinate tests. All the
variables are functions of $E$, so the deterministic parameter-law
identity \eqref{eq:deterministic-parameters} on $Z$ computes their
joint integral. At each fixed averaging time $t$, its state integral
is
\[
 \int_Z F_0(\omega_Z(S_Z^{bt}z))
    \prod_{i\in I}H_i\big(\kappa(\omega_Z(S_Z^{it}z))\big)
    \,\dd m_Z(z).
\]
Make the measure-preserving change of variable $z'=S_Z^{bt}z$.
The product of the $H_i$ becomes
\[
 \prod_{i\in I}H_i\big(
    \kappa(\omega_Z(S_Z^{(i-b)t}z'))\big).
\]
For each fixed $t$, this is measurable with respect to
$\kappa(\omega_Z(z'))$. To see this, the slot-$b$ map $Z\to Y$
is a flow factor, and the old factor $W_Y$ is preserved by every
flow time. Hence the old parameter at time $(i-b)t$ is a measurable
function of the old parameter at time zero.

Conditional expectation in this single-state integral therefore
replaces the first factor by
\[
 \overline F_0\bigl(\kappa(\omega_Z(z'))\bigr).
\]
Reversing the change of
variable and taking the deterministic averaging limit gives
\[
 \E_{\lambda_Z}\left[F_0(E_b)\prod_i H_i(p_i)\right]
 =\E_{\lambda_Z}\left[\overline F_0(w)\prod_i H_i(p_i)\right].
\]
The fixed-time identities suffice under the time integral by
measurability and Fubini; no simultaneous pointwise model for the
factor action is needed. The $E_b$-marginal is the single-state
parameter law, so $\overline F_0(w)=\E[F_0(E_b)\mid w]$.
Products of the $H_i$ determine the $p$-law, and $w=p_b$ is itself
a coordinate of $p$. The displayed identity is exactly conditional
independence given $w$.
\end{proof}

\subsection{Why the pairing cannot vanish}

Let $\alpha_w$ be the conditional law of the column parameter
array $p$ over its $b$th parameter $w$. Let $\beta_w$ be the
conditional law of $E_b$ over $w$. These are regular conditional
laws on standard spaces. Lemma~\ref{lem:parameter-independence}
gives
\begin{equation}
 \law(p,E_b\mid w)=\alpha_w\otimes\beta_w,
 \qquad R_*\beta_w=\alpha_w.
 \label{eq:parameter-disintegrations}
\end{equation}
The second identity uses the row marginal $m_Z=\lambda_Y$:
its old full parameter array $R(e)$ has the same law as the column
parameter array, with the same distinguished $b$th coordinate.

For almost every $w$, work in the separable real Hilbert space
\[
  \mathcal H_w=L^2(m_{Y,w};\R).
\]
The bounded functions $g(p,\cdot)$ and $u(e,\cdot)$ define
square-integrable random vectors in this space, with laws of their
parameters respectively $\alpha_w$ and $\beta_w$. There is no
measurability ambiguity in this fibrewise assertion. Bounded Borel
versions of $g$ and $u$ are fixed; disintegration and Fubini yield
their sections for almost every $w$. A countable generating algebra
of the standard Borel space $Y$ has rational simple span dense in
every $L^2(m_{Y,w})$. Tests against that span give weak measurability
of the sections, and separability gives strong measurability.

The conditional mean of the second random vector recovers the first:
\begin{equation}
 \E_{\beta_w}[u(e,\cdot)\mid R(e)]
     =g(R(e),\cdot)
 \quad\text{in }\mathcal H_w.
 \label{eq:fibre-tower}
\end{equation}
Here is the conditional-law justification. In a single $Z$-state,
write $y$ for its slot-$b$ state. Lemma~\ref{lem:row-column} gives,
conditional on $w$,
\[
  \law(e,y\mid w)=\beta_w(\dd e)m_{Y,w}(\dd y).
\]
Conditioning \eqref{eq:u} down on $(R(e),y)$ gives
$g(R(e),y)$ by \eqref{eq:q-and-g}, because $m_Z=\lambda_Y$.
Under the preceding product law, conditioning $e$ on $(R(e),y)$
is the same as conditioning it on $R(e)$. This proves
\eqref{eq:fibre-tower}. One can impose the scalar conditional
identities first for countably many generating tests and then use
Fubini and density to obtain the stated Hilbert-space identity for
almost every $w$.

Suppose that the right side of \eqref{eq:row-column} vanished
almost surely. By \eqref{eq:parameter-disintegrations}, for almost
every $w$ this would give
\[
  \langle g(p,\cdot),u(e,\cdot)\rangle_{\mathcal H_w}=0
  \quad(\alpha_w\otimes\beta_w)\text{-almost surely}.
\]
Condition on $R(e)=p'$ and use \eqref{eq:fibre-tower}. We obtain
\begin{equation}
 \langle g(p,\cdot),g(p',\cdot)\rangle_{\mathcal H_w}=0
 \quad(\alpha_w\otimes\alpha_w)\text{-almost surely}.
 \label{eq:independent-orthogonality}
\end{equation}
A square-integrable random vector in a separable real Hilbert
space cannot be orthogonal almost surely to an independent copy
unless it vanishes almost surely. Indeed, otherwise its norm is at
least some $\epsilon>0$ with positive probability. By a countable
cover with balls of radius $\epsilon/4$, this event has positive
probability in one such ball. Two independent copies then have
positive probability of both lying in that ball and having norm
at least $\epsilon$. Their distance is less than $\epsilon/2$, so
\[
 \langle v,v'\rangle
 =\tfrac12(\|v\|^2+\|v'\|^2-\|v-v'\|^2)>0,
\]
a contradiction. Apply this observation to
\eqref{eq:independent-orthogonality}. It forces $g(p,\cdot)=0$
for $\alpha_w$-almost every $p$, for almost every $w$.
The singleton conditional law for the column then implies
$g(P,y_b)=0$ in $L^2(\lambda_Y)$, contrary to
\eqref{eq:q-and-g}. Thus \eqref{eq:row-column} does not vanish
almost surely.

\subsection{Increasing the number of output roles}

Choose a bounded real function $\Theta(E)$ that detects the
nonzero conditional expectation in \eqref{eq:row-column}; its sign
is one such choice. Then
\begin{equation}
 \E_{\lambda_Z}\left[
   \Theta(E)q((y_i)_i)q(\mathbf z_b)\right]\ne0.
 \label{eq:improved-correlation}
\end{equation}
We interpret this as a correlation on the regrouped channel $Z$.
For $i\in J\setminus\{b\}$, the factor
$h_i(L_i(y_{ib}))$ retains its old role. An old output role uses
only the retained outputs from slot $b$. An old input role uses
the new original coordinate and retained outputs from its old
proper subset $H_i$ of group numbers. It may ignore the additional
outputs in those groups. Each such test is centered for its new
product spatial law, since its marginal on the coordinates it
actually uses is the old product law.

At row $b$, the factor $q(\mathbf z_b)$ reads only the added outputs
with spatial lists $Q_j$, $j\in J\setminus\{b\}$, and thus has an
output role. It is centered: under the product spatial law of all
new outputs these lists are independent product lists, so the mean
of this product of centered tests is zero. Hence
\eqref{eq:improved-correlation} is a nonzero correlation of support
size $k$ with $r+1$ output roles. This contradicts the maximal
choice of $r$.

\begin{proof}[Completion of the proof of Theorem~\ref{thm:main}]
If the required pointwise limit failed for a fixed tuple and a
fixed $n\geq2$, the reduction and selectors of
Section~\ref{sec:selected} would give the positive correlation
\eqref{eq:positive-witness}. The minimum-support argument rules this
out immediately when $n<4$, and the saturated, regrouped square
argument rules it out when $n\geq4$. Therefore the averages converge
almost everywhere to the product of the means, along all positive
integer lengths. Finally, the countable-coordinate coding reduction
pulls this conclusion back to the original, possibly nonstandard,
probability space. The tuple and $n$ were fixed throughout, as allowed
in the statement of the theorem.
\end{proof}

\appendix
\section{Tensor packing}\label{sec:packing}

The closure proof uses a deterministic fact: a subspace of moderately
growing dimension captures a substantial part of only a small fraction
of an almost orthogonal tensor family.  We reproduce the argument of
\cite[Theorem~4.1]{FourfoldCompanion}, with the same quantitative
estimate.  Throughout this appendix, Hilbert spaces are real, tensor
products are Hilbert tensor products, and logarithms are natural.

\begin{theorem}[Tensor packing]\label{thm:packing}
There is a universal constant $C$ with the following property.  Let
$0<\varepsilon\le1$, and let $d\ge1$ and $L\ge0$ be integers.
Let $x_1,\ldots,x_m$ and $y_1,\ldots,y_m$ lie in the unit balls of
real Hilbert spaces $H_1,H_2$.  For each family suppose that a graph
on $\{1,\ldots,m\}$ of maximum degree at most $L$ contains every
off-diagonal pair whose absolute inner product exceeds $\eta$, where
\[
 0\le\eta\le\varepsilon^2/10^4.
\]
The two exceptional graphs may differ.  For every orthogonal projection
$P$ on $H_1\otimes H_2$ of rank at most $d$,
\begin{equation}\label{eq:packing-quantitative}
 \#\{j:\norm{P(x_j\otimes y_j)}\ge\varepsilon\}
 \le C(1+L)\varepsilon^{-2}
       \exp\!\left(C\varepsilon^{-2}\sqrt d\log(2+d)\right).
\end{equation}
For $d=d_m$ as defined in \eqref{eq:channel-dimension}, this is
$o_{\varepsilon,L}(m)$, uniformly over the Hilbert spaces, vectors,
projections, and allowed values of $\eta$.
\end{theorem}

The proof separates directions on which the testing space has large
variance and applies Gaussian comparison to the remaining directions.
The comparison is the finite Sudakov--Fernique inequality; see
\cite[Section~1, Equation~(1.2)]{Fernique1974} and
\cite{Sudakov1971}.  The following matrix estimate uses the Gaussian
integration-by-parts moment argument behind matrix Khintchine
inequalities; compare \cite[Lemma~6.1 and Proposition~7.3]{Tropp2018}.
We include its proof to retain dependence on the covariance trace,
rather than on the ambient dimensions.

\begin{lemma}[Gaussian matrix estimate]\label{lem:packing-gaussian}
Let $B_1,\ldots,B_r$ be matrices between two finite-dimensional real
Hilbert spaces, and put
\[
 H_j=\begin{pmatrix}0&B_j\\B_j^*&0\end{pmatrix},
 \qquad S=\sum_{j=1}^rH_j^2.
\]
If $\norm{S}_{\rm op}\le\sqrt d$ and $\tr S\le2d$ for some
$d\ge1$, then for independent standard normal variables $g_j$,
\begin{equation}\label{eq:packing-gaussian-bound}
 \E\left\|\sum_{j=1}^rg_jB_j\right\|_{\rm op}
       \le C d^{1/4}\sqrt{\log(2+d)}.
\end{equation}
The constant is independent of the dimensions of the two spaces.
\end{lemma}

\begin{proof}
Write $H=\sum_jg_jH_j$.  Gaussian integration by parts, followed by
differentiation of the matrix power, gives for every integer $q\ge2$
\[
 \E\tr H^{2q}
 =\sum_{j=1}^r\sum_{a=0}^{2q-2}
      \E\tr(H_jH^aH_jH^{2q-2-a}).
\]
For a fixed self-adjoint $H$ with eigenvalues $\lambda_u$ and a
self-adjoint matrix $K$, the absolute value of such a trace is at most
\[
 \sum_{u,v}|K_{uv}|^2
       |\lambda_v|^a|\lambda_u|^{2q-2-a}
 \le\tr(K^2|H|^{2q-2}).
\]
The inequality follows from weighted arithmetic--geometric mean and
the symmetry $|K_{uv}|=|K_{vu}|$.  Summing over $j$ gives
\begin{equation}\label{eq:packing-moment-recursion}
 \E\tr H^{2q}
 \le(2q-1)\norm{S}_{\rm op}\E\tr|H|^{2q-2}.
\end{equation}
Since $\E\tr H^2=\tr S$, iteration yields
\[
 \E\tr H^{2q}
 \le(2q-1)!!\,\norm{S}_{\rm op}^{q-1}\tr S.
\]
The operator norm of $H$ equals that of $\sum_jg_jB_j$.
Taking a $2q$-th root and using the hypotheses bounds its expectation by
\[
 C\sqrt q\,d^{1/4}(2\sqrt d)^{1/(2q)}.
\]
Choose an integer $q\ge2$ comparable to $\log(2+d)$ to obtain
\eqref{eq:packing-gaussian-bound}.
\end{proof}

\begin{proof}[Proof of Theorem~\ref{thm:packing}]
We first record an elementary consequence of the exceptional-graph
hypothesis.  If unit-ball vectors $v_j$ have absolute inner products
at most $\eta$ outside a graph of degree $L$, then for every nonempty
set $A$ of indices,
\begin{equation}\label{eq:packing-cluster}
 \left\|\frac1{|A|}\sum_{j\in A}v_j\right\|^2
       \le\frac{1+L}{|A|}+\eta.
\end{equation}
Indeed the expansion contains at most $(1+L)|A|$ ordered diagonal or
exceptional pairs.  A lower bound on the norm of an average therefore
bounds the number of vectors in that average.

It suffices to work in the finite spans of the $x_j$ and $y_j$.
To see this, let $Q$ project onto their tensor product and replace the
testing subspace $\operatorname{Ran}P$ by $Q(\operatorname{Ran}P)$.
For any vector $z$ in that tensor product, the norm of its projection
onto the new subspace is at least $\norm{Pz}$: in the supremum defining
$\norm{Pz}$, replace each unit testing vector by its $Q$-image, whose
norm is at most one and whose pairing with $z$ is unchanged.  The
dimension does not increase.

Identify tensors with matrices from $H_2$ to $H_1$.  Choose a
Hilbert--Schmidt orthonormal basis $A_1,\ldots,A_r$ of the testing
space, where $r\le d$.  The matrices
\[
 R_1=\sum_jA_jA_j^*,\qquad R_2=\sum_jA_j^*A_j
\]
have trace $r$.  Let $E_1,E_2$ be their spectral subspaces for
eigenvalues strictly greater than $\sqrt d$; each has dimension at most
$\sqrt d$.

Discard indices for which the projection of $x_j$ onto $E_1$ has norm
greater than $\varepsilon/10$, and do the same for $y_j,E_2$.
Cover either exceptional subspace's unit ball by at most
$(1+80/\varepsilon)^{\sqrt d}$ balls of radius $\varepsilon/40$.
If one ball contains $s$ of the discarded projections, their average
has norm at least $\varepsilon/20$.  Applying
\eqref{eq:packing-cluster} to the original vectors, with
$\eta\le\varepsilon^2/10^4$, gives
$s\le C(1+L)\varepsilon^{-2}$.  The number of discarded indices is
therefore at most
\begin{equation}\label{eq:packing-discarded}
 C(1+L)\varepsilon^{-2}
   \exp\!\bigl(\sqrt d\log(1+80/\varepsilon)\bigr).
\end{equation}

Let $Q_1,Q_2$ project onto $E_1^\perp,E_2^\perp$, and put
$B_j=Q_1A_jQ_2$.  The coefficient map
$z\mapsto(\ip{A_j}{z}_{\rm HS})_{j=1}^r$ is a contraction.
For a remaining index,
\[
 \norm{x_j\otimes y_j-Q_1x_j\otimes Q_2y_j}
 \le\norm{(1-Q_1)x_j}+\norm{(1-Q_2)y_j}
 \le\varepsilon/5.
\]
Thus each remaining index counted in
\eqref{eq:packing-quantitative} has coefficient vector
\[
 v_j=(\ip{x_j}{B_ay_j})_{a=1}^r,
 \qquad\norm{v_j}\ge\varepsilon/2.
\]
The map from $x_j\otimes y_j$ to $v_j$ remains contractive.
Outside the union of the two exceptional graphs, the tensor inner
products have absolute value at most $\eta^2$, and that union has
degree at most $2L$.  A collection of these coefficient vectors in a
ball of radius $\varepsilon/8$ centered at one of them has average
norm at least $3\varepsilon/8$.  Applying
\eqref{eq:packing-cluster} to its preimage tensors bounds its size by
$C(1+L)\varepsilon^{-2}$.  A maximal $\varepsilon/8$-separated subset
therefore has cardinality $l$ such that the number of indices still
to be counted is at most $C(1+L)\varepsilon^{-2}l$.

We now bound this separated set.  For $G=\sum_ag_aB_a$, the centered
Gaussian variables $Z_j=\ip{x_j}{Gy_j}$ satisfy
\[
 \E|Z_j-Z_{j'}|^2=\norm{v_j-v_{j'}}^2.
\]
Gaussian comparison gives, when $l\ge2$,
\begin{equation}\label{eq:packing-comparison}
 c\varepsilon\sqrt{\log l}
       \le\E\max_j Z_j\le\E\norm{G}_{\rm op}.
\end{equation}
For completeness, the comparison of finite centered Gaussian families
follows by applying Gaussian integration by parts to
$\beta^{-1}\log\sum_j e^{\beta z_j}$ along their independent Gaussian
interpolation.  Its derivative is
\[
 \frac\beta4\E\sum_{j,j'}p_jp_{j'}
 \left(\E|Z_j-Z_{j'}|^2-\E|Z'_j-Z'_{j'}|^2\right),
 \qquad p_j=\frac{e^{\beta z_j}}{\sum_a e^{\beta z_a}}.
\]
It is nonnegative when the first increment variances dominate.
Letting $\beta\to\infty$ compares the expected maxima.  For
\eqref{eq:packing-comparison}, take the $Z'_j$ to be independent
normals of variance $\varepsilon^2/128$ and use
$\E\max_{j\le l}g_j\ge c\sqrt{\log l}$.

Compression gives
\[
 \sum_jB_jB_j^*\le Q_1R_1Q_1,\qquad
 \sum_jB_j^*B_j\le Q_2R_2Q_2.
\]
These operators have norms at most $\sqrt d$ and traces at most $d$.
Lemma~\ref{lem:packing-gaussian} therefore applies.  Together with
\eqref{eq:packing-comparison} it yields
\[
 \log l\le C\varepsilon^{-2}\sqrt d\log(2+d).
\]
The resulting bound holds for $l\le1$ as well.  Combine it with
\eqref{eq:packing-discarded} and absorb
$\log(1+80/\varepsilon)\le C\varepsilon^{-2}$ to prove
\eqref{eq:packing-quantitative}.

For $d=d_m$, the exponent is
$O_\varepsilon((\log m)^{3/4}\log\log m)=o(\log m)$,
which proves the asserted $o(m)$ bound.
\end{proof}

Since projection norms are at most one, the estimate also gives
\begin{equation}\label{eq:packing-average}
 \frac1m\sum_{j=1}^m\norm{P(x_j\otimes y_j)}
 \le\varepsilon+
 \frac{C(1+L)}{m\varepsilon^2}
 \exp\!\left(C\varepsilon^{-2}\sqrt{d_m}\log(2+d_m)\right)
\end{equation}
when $\operatorname{rank}P\le d_m$.  Here $\varepsilon$ and then
$\eta$ are fixed before $m$ increases; $L$ may depend on that chosen
$\eta$.  This is exactly the order of quantifiers supplied by mixing
in Proposition~\ref{prop:channel-closure}.

\bibliographystyle{amsplain}
\begingroup\raggedright
\bibliography{references}
\endgroup
\end{document}
